
\documentclass{article}
\usepackage{colortbl}
\usepackage{makecell}
\usepackage{multirow}
\usepackage{supertabular}

\begin{document}

\newcounter{utterance}

\centering \large Interaction Transcript for game `referencegame', experiment `number\_grids', episode 4 with reflexion{-}Qwen3.5{-}9B.
\vspace{24pt}

{ \footnotesize  \setcounter{utterance}{1}
\setlength{\tabcolsep}{0pt}
\begin{supertabular}{c@{$\;$}|p{.15\linewidth}@{}p{.15\linewidth}p{.15\linewidth}p{.15\linewidth}p{.15\linewidth}p{.15\linewidth}}
    \# & $\;$A & \multicolumn{4}{c}{Game Master} & $\;\:$B\\
    \hline

    \theutterance \stepcounter{utterance}  
    & & \multicolumn{4}{p{0.6\linewidth}}{
        \cellcolor[rgb]{0.9,0.9,0.9}{
            \makecell[{{p{\linewidth}}}]{
                \texttt{\tiny{[P1$\langle$GM]}}
                \texttt{You have attempted this same task before and did not succeed. These are the plans you wrote after those attempts.} \\
\texttt{Trial \#1: New plan: My previous description was ambiguous because it only specified the fourth row and implied the rest of the grid's structure without explicitly defining the pattern for the other rows, which could match the distractors. Specifically, Distractor 1 has a different fourth row, but Distractor 2 has a different pattern in the third and fourth rows. To ensure uniqueness, I must explicitly describe the alternating pattern of the filled cells in the odd{-}numbered rows (rows 1, 3, 5) versus the even{-}numbered rows (rows 2, 4). The correct strategy is to state that rows 1, 3, and 5 are fully filled except for the first and fourth columns, while rows 2 and 4 have empty cells in the third and fourth columns. I will generate an expression that clearly distinguishes the target's specific row{-}by{-}row configuration from both distractors by mentioning the exact column indices of the filled cells for the repeating pattern.} \\
\\ 
\texttt{Trial \#2: New plan: My previous description failed because it relied on an "alternating row pattern" that was actually shared by Distractor 2, making the expression non{-}unique. Specifically, both the target and Distractor 2 have identical patterns in rows 1, 2, 3, and 5. The only distinguishing feature is the fourth row: in the target grid, the fourth row has filled cells in columns 2 and 5 (matching the vertical bar), whereas in Distractor 2, the fourth row is completely empty. To succeed, I must explicitly describe the fourth row's specific configuration ("the fourth row has filled cells only in the second and fifth columns") rather than relying on a general alternating pattern, ensuring the description matches the target but excludes both distractors.} \\
\\ 
\\ 
\texttt{You are given three grids, where each of them is 5 by 5 in size.} \\
\texttt{Grids have empty cells marked with "▢" and filled cells marked with "X".} \\
\texttt{Your task is to generate a referring expression that best describes the target grid while distinguishing it from the two other distractor grids.} \\
\texttt{The first grid is the target grid, and the following two grids are the distractors.} \\
\\ 
\texttt{Target grid:} \\
\\ 
\texttt{□ X X X X} \\
\texttt{□ X □ □ X} \\
\texttt{□ X X X X} \\
\texttt{□ X □ □ X} \\
\texttt{□ X X X X} \\
\\ 
\texttt{Distractor grid 1:} \\
\\ 
\texttt{□ X X X X} \\
\texttt{□ X □ □ X} \\
\texttt{□ X X X X} \\
\texttt{□ □ □ □ □} \\
\texttt{□ X X X X} \\
\\ 
\texttt{Distractor grid 2:} \\
\\ 
\texttt{□ X X X X} \\
\texttt{□ X □ □ X} \\
\texttt{□ X □ □ X} \\
\texttt{□ X □ □ X} \\
\texttt{□ X X X X} \\
\\ 
\texttt{Instruction: Describe the target grid.} \\
\texttt{Generate the referring expression starting with the tag "Expression: " for the given target grid. Omit any other text.} \\
            }
        }
    }
    & \\ \\

    \theutterance \stepcounter{utterance}  
    & \multicolumn{4}{p{0.6\linewidth}}{
        \cellcolor[rgb]{0.9,0.9,0.9}{
            \makecell[{{p{\linewidth}}}]{
                \texttt{\tiny{[P1$\rangle$GM]}}
                \texttt{Expression: The grid has a vertical column of filled cells in the second column spanning all five rows. Additionally, rows 1, 3, and 5 are fully filled except for the first and fourth columns, while rows 2 and 4 have filled cells only in the second and fifth columns.} \\
            }
        }
    }
    & & \\ \\

    \theutterance \stepcounter{utterance}  
    & & & \multicolumn{2}{p{0.3\linewidth}}{
        \cellcolor[rgb]{0.9,0.9,0.9}{
            \makecell[{{p{\linewidth}}}]{
                \texttt{\tiny{[GM$|$GM]}}
                \texttt{Expression: The grid has a vertical column of filled cells in the second column spanning all five rows. Additionally, rows 1, 3, and 5 are fully filled except for the first and fourth columns, while rows 2 and 4 have filled cells only in the second and fifth columns.} \\
            }
        }
    }
    & & \\ \\

    \theutterance \stepcounter{utterance}  
    & & \multicolumn{4}{p{0.6\linewidth}}{
        \cellcolor[rgb]{0.9,0.9,0.9}{
            \makecell[{{p{\linewidth}}}]{
                \texttt{\tiny{[GM$\rangle$P2]}}
                \texttt{You have attempted this same task before and did not succeed. These are the plans you wrote after those attempts.} \\
\texttt{Trial \#1: New plan: I incorrectly identified the First grid because I overlooked the specific pattern in the fourth row. The expression requires the fourth row to be entirely empty ("▢"), whereas the First grid has an "X" in the second column of that row. The Second and Third grids both have an "X" in the fourth row's second column, but the Third grid is the only one where the fourth row is completely empty (rows 1{-}3 and 5 have Xs in cols 2,3,5; row 4 is all empty). Next time, I will carefully verify the fourth row of each grid against the "only empty cells" constraint before making a selection. The correct answer is Third.} \\
\\ 
\texttt{Trial \#2: New plan: I incorrectly selected the First grid because I failed to rigorously verify the "alternating row pattern" for the remaining columns (3, 4, 5) against the specific expression rules. The expression mandates that in even{-}numbered rows (2 and 4), only the fifth column is filled, while the First grid has filled cells in columns 3 and 4 for row 2. I will now systematically cross{-}reference every cell of the Second and Third grids against the expression's two main constraints: 1) Column 1 is empty and Column 2 is full, and 2) Rows 1, 3, 5 have Xs in cols 3{-}5, while Rows 2, 4 have Xs only in col 5. Upon this verification, the Third grid is the only one that satisfies the condition where Row 4 has only the fifth column filled (matching the pattern), whereas the Second grid has an extra X in the fourth column of Row 4. The correct answer is Third.} \\
\\ 
\\ 
\texttt{You are given three grids, where each of them is 5 by 5 in size.} \\
\texttt{Grids have empty cells marked with "▢" and filled cells marked with "X".} \\
\texttt{You are also given a referring expression that describes one of the given grids.} \\
\texttt{Your task is to select the grid that matches the given referring expression.} \\
\texttt{Generate only the number (in text) of the grid that matches the given expression by selecting first, second, or third.} \\
\\ 
\texttt{First:} \\
\\ 
\texttt{□ X X X X} \\
\texttt{□ X □ □ X} \\
\texttt{□ X X X X} \\
\texttt{□ □ □ □ □} \\
\texttt{□ X X X X} \\
\\ 
\texttt{Second:} \\
\\ 
\texttt{□ X X X X} \\
\texttt{□ X □ □ X} \\
\texttt{□ X X X X} \\
\texttt{□ X □ □ X} \\
\texttt{□ X X X X} \\
\\ 
\texttt{Third:} \\
\\ 
\texttt{□ X X X X} \\
\texttt{□ X □ □ X} \\
\texttt{□ X □ □ X} \\
\texttt{□ X □ □ X} \\
\texttt{□ X X X X} \\
\\ 
\texttt{Expression: The grid has a vertical column of filled cells in the second column spanning all five rows. Additionally, rows 1, 3, and 5 are fully filled except for the first and fourth columns, while rows 2 and 4 have filled cells only in the second and fifth columns.} \\
\texttt{Question: Which grid does the expression refer to?} \\
\texttt{Start with the tag "Answer: ", followed by your selection. Omit any other text.} \\
            }
        }
    }
    & \\ \\

    \theutterance \stepcounter{utterance}  
    & & & \multicolumn{4}{p{0.6\linewidth}}{
        \cellcolor[rgb]{0.9,0.9,0.9}{
            \makecell[{{p{\linewidth}}}]{
                \texttt{\tiny{[GM$\langle$P2]}}
                \texttt{Answer: Third} \\
            }
        }
    }
     \\ \\

    \theutterance \stepcounter{utterance}  
    & & & \multicolumn{2}{p{0.3\linewidth}}{
        \cellcolor[rgb]{0.9,0.9,0.9}{
            \makecell[{{p{\linewidth}}}]{
                \texttt{\tiny{[GM$|$GM]}}
                \texttt{third} \\
            }
        }
    }
    & & \\ \\

    \theutterance \stepcounter{utterance}  
    & \multicolumn{4}{p{0.6\linewidth}}{
        \cellcolor[rgb]{0.95,0.95,0.95}{
            \makecell[{{p{\linewidth}}}]{
                \texttt{\tiny{[P2$\rightarrow$P1]}}
                \texttt{Post{-}episode Reflexion (not sent to GM):} \\
\texttt{New plan: My previous attempts failed because they relied on describing the repeating pattern of rows 1, 3, and 5, which is identical in both the target and Distractor 2, making the description non{-}unique. The only feature that distinguishes the target from both distractors is the specific configuration of the fourth row: in the target, it contains filled cells in columns 2 and 5 (matching the vertical bar), whereas in Distractor 1 it is completely empty, and in Distractor 2 it is also completely empty. To succeed, I must construct an expression that explicitly defines the entire grid by stating the unique pattern of the fourth row ("the fourth row has filled cells only in the second and fifth columns") while confirming the consistent vertical bar in the second column and the specific filled cells in the odd{-}numbered rows, ensuring no ambiguity remains that could apply to the distractors.} \\
            }
        }
    }
    & & \\ \\

    \theutterance \stepcounter{utterance}  
    & & & \multicolumn{4}{p{0.6\linewidth}}{
        \cellcolor[rgb]{0.95,0.95,0.95}{
            \makecell[{{p{\linewidth}}}]{
                \texttt{\tiny{[P1$\rightarrow$P2]}}
                \texttt{Post{-}episode Reflexion (not sent to GM):} \\
\texttt{New plan: I failed because I did not rigorously verify the specific pattern of filled cells in the even{-}numbered rows (rows 2 and 4) against the expression's requirement that they contain filled cells *only* in the second and fifth columns. While I correctly identified that the First grid was invalid due to its fourth row, I likely overlooked the discrepancy in the Second grid where row 4 has an extra filled cell in the fourth column, violating the "only" constraint. Next time, I will systematically check every grid against the two critical constraints: 1) Column 2 must be fully filled across all rows, and 2) Rows 2 and 4 must have filled cells exclusively in columns 2 and 5, with all other cells in those rows being empty. Upon this strict verification, the Third grid is the only one that satisfies the condition where Row 4 has only the fifth column filled (along with the second), whereas the Second grid fails this specific check. The correct answer is Third.} \\
            }
        }
    }
     \\ \\

\end{supertabular}
}

\end{document}
